\newcommand{\PrimaryRuns}{840}
\newcommand{\ContrastCoverageRange}{10.4\% to 90.0\%}
\newcommand{\RawFlipRange}{16.0--39.3\%}
\newcommand{\ResultsCoverage}{The selected contrast certifies 81.8\% coverage on CR6261--H1, 77.9\% on CH65--MA90, and 77.3\% on CH65--G189E. It reaches only 10.4\% on CR9114--H3 and 17.4\% on CR6261--H9. The published full-curve fits reach 92.5\% and 28.9\% on the first and fourth tasks, respectively. The sparse rule therefore offers a useful tradeoff in some settings and a substantial loss in others. No primary run has aggregate accepted error above 10\% (0/840 violations); empty acceptance sets are excluded from error averages. This observed statement does not replace the probabilistic guarantee or establish simultaneous coverage for all runs. All eight methods defer on the influenza-B cohorts, which contain only 115--149 evaluable variants. Their calibration subsets contain at most 44 variants before class-specific acceptance, leaving insufficient evidence under the declared correction.}
\newcommand{\ResultsRecall}{On CH65--SI06, the fixed three-point rule covers 98.7\% of the evaluable test cohort with 2.1\% accepted error, but retains only 29.8\% of reference hits for further testing. The raw full-panel learner retains 95.7\% at 92.6\% coverage. No method positively identifies any test hit on CH65--SI06 or CR9114--H1 after class-specific calibration; its retained hits, where present, are deferred rather than declared discoveries. The hit-protection variant makes this distinction operational. With the same 300-variant calibration cap, every method defers completely on CH65--SI06. It has no observed class-error or hit-loss violation in its 504 method--split runs, but often achieves this by sharply reducing negative acceptance. Appendix tables report positively identified hits and hits retained separately.}
\newcommand{\ResultsStress}{The constructed test-only transformations change 16.0--39.3\% of forced decisions across raw-input method/task means. Both the learned three-point ratio and the full-panel invariant learner have zero decision changes in all reported runs. The standardized contrast also preserves its decisions and acceptance set when its linear-space spreads transform consistently. Raw learners can violate their original accepted-error target after the shift; this is an expected failure of transporting a calibration certificate outside its conditions, not a counterexample to the theorem.}
\newcommand{\ResultsSensitivity}{Changing the calibration cap from 100 to 1,000 variants changes which gates can be certified, without changing the physical number of source conditions. Changing the cutoff by half a log unit changes both reference prevalence and the comparison being asked. The supplementary tables include every declared sensitivity setting, including 5\% and 20\% error targets; no setting is omitted because its coverage is zero. The retained-hit analysis additionally shows that a 10\% accepted-error target and a 10\% discarded-hit target are materially different objectives.}
\newcommand{\VerificationText}{Independent data alignment reconstructs stored fit $r^2$ for 18,475 high-quality source curves across all 21 tables; the maximum absolute discrepancy is 5.45e-10. The verification suite checks 55,000 noiseless contrast signs under varied affinities, amplitudes, backgrounds, and affine transforms. It exhaustively checks exact bound coverage in 6,273 finite-population settings, and verifies that changing held-out test labels leaves decisions and calibration gates unchanged. These checks establish implementation properties under their stated inputs, not wet-lab validity of the response model.}
